% H44: erasure, re-acquisition and thrash. Sources: hypotheses/H44-erasure-reacquisition-thrash/{README.md, summary/content.tex, summary/meta.json};
% writeup/visuals/H44-context-wipe/{caption.tex, README.md}; interpretation/swarm-constants.json (ΔW_erase).
\section{H44: A context wipe costs output but does not heat the agent}

\textbf{The question:}
The scaffold wipes an agent's context window at a cap of 41 records (40 model calls) and keeps its memory.
Does the agent then thrash: re-read instead of produce, act more randomly and listen harder to the room?
How much output does each wipe cost?

\textbf{The intuition:}
Treat each model call as a spin with many states (write, run, read, look, talk, idle).
Three fields set its choice: a baseline field (what the agent does with no task in mind), a room field and a task field.
The context window holds the task field.
The wipe is a quench, a sudden change of a control parameter: here it sets the task field to zero.
A \emph{field quench} shifts the mix of calls toward reading, and the task field then grows back over a few calls.
A \emph{temperature pulse} makes every choice more random, so the entropy of the call mix rises.
Linear response, the rule that a spin with a weak self-field follows an outside field more easily, adds one prediction.
An agent with no task field should follow the room more closely after a wipe.

\textbf{What we measured:}
The unit is one goal period: nine regime-III periods outside the reserved data (G36--G51), with 1.24M calls, 21{,}165 forced and 16{,}357 voluntary resets.
The thrash index $\Theta_c$ is the rise in the read share of non-write calls over calls $+1\ldots+5$.
It conditions on the agent and on the category of the previous call.
The write dip $\Omega$ compares writes over calls $+1\ldots+10$ with calls $-20\ldots-11$; DQ4 work commits give a second output measure.
The null is a set of pseudo-erasures: no-reset boundaries at call 21 or 31 of the same segments.
They match the hour of day by construction (KS distance $\le0.04$).
\emph{Scheduler field:} removed, because each event compares two sides of one reset within one agent-day.
\emph{External drive:} removed, because the cap, not the task state, times a forced wipe.
The forced pre-trend ($+7\%$) equals the no-reset band.
\emph{Shared model prior:} removed by the conditioning on agent and previous call; the effect holds in Anthropic, OpenAI and Google agents.
\emph{Common input:} the reply test is a difference-in-differences against pseudo-erasures, so a common response cancels.
\emph{Synthetic check:} three planted worlds (thrash, pure output dip, no effect) on real call skeletons.
At G51 counts the pipeline classifies all three correctly in 100\% of replicates.
The check also showed that the raw index calls a pure dip ``thrash''; hence the conditioning in $\Theta_c$.

\begin{figure}[h]
\includegraphics[width=\textwidth]{H44-context-wipe/fig.pdf}
\caption{What a forced context wipe does to one agent.
In G51 the share of re-reading calls doubles at the first call after the wipe. Writes and work commits dip for about three calls; the gray band is the no-reset null.
The lower panels show the work-commit dip in every period and the drop in command loops after a wipe.
Animation: \texttt{writeup/visuals/H44-context-wipe/anim.mp4}.}
\label{fig:int-H44}
\end{figure}

\textbf{What we found:}
Re-acquisition is real; thrash is not.
In G51 the read share rises from 0.22 to 0.47 at the first call and relaxes over 5--7 calls.
$\Theta_c=+0.079$ [0.071, 0.086], positive in 9/9 periods; the null gives $-0.011$.
Writes fall 26\% [$-29$, $-23$] (8/9 periods) and work commits fall 38\% [$-41$, $-35$].
Real command failures rise by 1 percentage point.
Per period the wipes cost 3.7--8.7\% of all write calls and 4.8--10.9\% of all work commits.
The entropy of the call mix \emph{falls} ($-0.027$): the agent reads in an orderly way, with no temperature pulse.
The reply rate to new room items does not change (rate ratio 0.97 [0.91, 1.04]; G51 0.90).
Coupling to items read before the wipe falls (ratio 0.60 [0.44, 0.81]).
A command loop recurs after a wipe in 16\% of cases, against 72\% without one (G38; pooled odds ratio 0.11 [0.05, 0.26], 7 periods).
Scope: regime III, nine periods.
The claim ``a forced wipe triggers orderly re-reading, costs 4--11\% of output and breaks loops'' stands\cred{0.88}.
Its expected usefulness is EU~3.23 (value if true 3.67; three raters, one blind).
The card and \texttt{meta.json} agree.

\textbf{What it means:}
For a scaffold builder, the cap length is a cost/benefit dial.
A 40-call cap costs about 4--11\% of output and breaks most command loops (odds ratio about 0.1).
Point a recovering agent at its working files.
An agent that first re-reads its files resumes writing 2.6 calls sooner than one that first looks at the screen.
Do not count on the memory note written at the wipe: its length has no effect on the dip ($\rho=-0.03$).
Do not expect a freshly wiped agent to be easier for its peers to steer.

\textbf{Caveats:}
A regex classifier assigns the call categories, and nobody has checked it against hand labels.
GUI browsing cannot be separated from production.
Writes ramp up through each 40-call segment, so the size of the dip depends on the reference window.
The re-acquisition paths are the agent's choice, so the ``files first'' result is observational.
The analysis runs about 12 statistics on nine periods with no correction; G37 has low power.
The fragile flag is not set.
No confirmation test on reserved goal periods has run: \texttt{confirm.py} has only run on stand-ins.
